\documentclass[12pt,a4paper]{article}

\usepackage[T2A]{fontenc}
\usepackage[utf8]{inputenc}
\usepackage[russian]{babel}
\usepackage{amsmath,amssymb}
\usepackage{enumitem}
\usepackage{geometry}
\usepackage{hyperref}

\geometry{margin=2cm}
\setlist{nosep}
\hypersetup{
    colorlinks=true,
    linkcolor=blue,
    urlcolor=blue
}

\DeclareMathOperator{\End}{End}
\DeclareMathOperator{\Ker}{ker}
\DeclareMathOperator{\Image}{im}
\DeclareMathOperator{\rank}{rank}
\DeclareMathOperator{\GL}{GL}
\DeclareMathOperator{\Mat}{M}
\DeclareMathOperator{\tr}{tr}

\newcommand{\Lin}[1]{\left\langle #1 \right\rangle}

\begin{document}

\begin{center}
    {\Large \textbf{Глава 3. Линейные операторы}}\\[4pt]
    {\large Краткий конспект по экзаменационным билетам 22--37}
\end{center}

\tableofcontents
\newpage

\section*{Что брать из каждого параграфа главы 3}
\addcontentsline{toc}{section}{Что брать из каждого параграфа главы 3}

\begin{itemize}
    \item \textbf{\S 1. Определения и примеры}: темы 22, 26, 27, 28, 29, 30 частично. Нужны: определение линейного оператора, примеры, операции над операторами, $\Image L$, $\Ker L$, определение алгебры, обратимый оператор, теорема о существовании оператора по образам базиса, критерии обратимости.
    \item \textbf{\S 2. Матрица линейного оператора}: темы 24, 26, 27, 28, 29. Нужны: построение матрицы по образам базисных векторов, формула $Y=AX$, ранг и дефект, теорема $r(L)+d(L)=n$, изоморфизм операторов и матриц, изоморфизм групп обратимых операторов и матриц.
    \item \textbf{\S 3. Закон изменения матрицы при изменении базиса}: темы 23, 25. Нужны: матрица перехода, формулы координат $X=CX'$, $X'=C^{-1}X$, формула подобия $B=C^{-1}AC$.
    \item \textbf{\S 4. Собственные векторы}: темы 31, 32, 33. Нужны: собственные векторы и значения, собственные подпространства, линейная независимость собственных векторов, характеристический многочлен, спектр, алгоритмы нахождения собственных значений и векторов.
    \item \textbf{\S 5. Диагонализируемые операторы}: темы 34, 35. Нужны: определение диагонализируемости, критерий через базис из собственных векторов, критерий через кратности характеристических корней.
    \item \textbf{\S 6. Инвариантные подпространства}: темы 36, 37. Нужны: определение, примеры, связь с блочно-треугольными и блочно-диагональными матрицами, индуцированный оператор, делимость характеристических многочленов.
\end{itemize}

\textbf{Можно пропустить, если преподаватель не требует дополнительно:} биографические вставки, олимпиадные упражнения, минимальный многочлен, корневые подпространства, идемпотентные операторы, разложение Фиттинга и полный параграф про триангуляризуемые операторы.

\section{Билет 22. Операторы линейных пространств}

\subsection*{Определение}

Пусть $V$ и $W$ --- линейные пространства над полем $P$. Отображение
\[
    L:V\to W
\]
называется \textbf{линейным оператором}, если выполняются условия:
\[
    L(x+y)=L(x)+L(y),
\]
\[
    L(\alpha x)=\alpha L(x).
\]
Эквивалентно эти два условия можно записать одним:
\[
    L(\alpha x+\beta y)=\alpha L(x)+\beta L(y).
\]

\subsection*{Примеры}

\begin{itemize}
    \item тождественный оператор: $I(x)=x$;
    \item нулевой оператор: $O(x)=0$;
    \item оператор подобия: $(\alpha I)(x)=\alpha x$;
    \item поворот плоскости;
    \item проекция, например $L:\mathbb{R}^{3}\to\mathbb{R}^{2}$, $L(x,y,z)=(x,y)$;
    \item дифференцирование многочленов;
    \item интегрирование многочленов с фиксированной константой.
\end{itemize}

\subsection*{Теорема о задании оператора на базисе}

Пусть $a_1,\ldots,a_n$ --- базис $V$, а $b_1,\ldots,b_n$ --- произвольные векторы из $W$. Тогда существует единственный линейный оператор $L:V\to W$, такой что
\[
    L(a_i)=b_i,\qquad i=1,\ldots,n.
\]

\subsection*{Идея доказательства}

Любой $x\in V$ единственно представим в виде
\[
    x=x_1a_1+\cdots+x_na_n.
\]
Поэтому оператор обязан действовать так:
\[
    L(x)=x_1b_1+\cdots+x_nb_n.
\]
Это задает оператор однозначно; линейность проверяется прямой подстановкой.

\section{Билет 23. Связь между двумя базисами и координатами}

Пусть $e_1,\ldots,e_n$ --- старый базис, а $f_1,\ldots,f_n$ --- новый базис. Каждый новый базисный вектор раскладывается по старому:
\[
    f_j=c_{1j}e_1+\cdots+c_{nj}e_n.
\]

Матрица $C=(c_{ij})$ называется \textbf{матрицей перехода от старого базиса к новому}. Ее столбцы --- координаты новых базисных векторов в старом базисе:
\[
    (f_1,\ldots,f_n)=(e_1,\ldots,e_n)C.
\]

Если $X$ --- координаты вектора $x$ в старом базисе, а $X'$ --- координаты того же вектора в новом базисе, то
\[
    X=CX',
\]
\[
    X'=C^{-1}X.
\]

Главное: чтобы перейти от старых координат к новым, нужно умножить на $C^{-1}$.

\section{Билет 24. Матрица оператора}

Пусть $L:V\to W$, $a_1,\ldots,a_n$ --- базис $V$, $b_1,\ldots,b_m$ --- базис $W$.

Раскладываем образы базисных векторов:
\[
    L(a_i)=\alpha_{1i}b_1+\cdots+\alpha_{mi}b_m.
\]

Матрица
\[
    A_L=(\alpha_{ji})
\]
называется \textbf{матрицей оператора} $L$ в выбранных базисах.

\textbf{Важно:} в $i$-м столбце матрицы стоят координаты вектора $L(a_i)$.

Если
\[
    x=x_1a_1+\cdots+x_na_n,
\]
\[
    L(x)=y_1b_1+\cdots+y_mb_m,
\]
то в координатах
\[
    Y=AX.
\]

То есть матрица оператора переводит столбец координат $x$ в столбец координат $L(x)$.

\section{Билет 25. Матрицы оператора в разных базисах}

Пусть $L\in\End V$. В старом базисе матрица оператора равна $A$, в новом --- $B$. Пусть $C$ --- матрица перехода от старого базиса к новому. Тогда
\[
    B=C^{-1}AC.
\]

\subsection*{Короткое доказательство}

Для координат одного вектора:
\[
    X=CX'.
\]
В старом базисе действие оператора:
\[
    Y=AX.
\]
В новом:
\[
    Y'=BX',
\]
а также
\[
    Y=CY'.
\]
Тогда
\[
    CY'=ACX',
\]
откуда
\[
    Y'=C^{-1}ACX'.
\]
Следовательно,
\[
    B=C^{-1}AC.
\]

\subsection*{Подобные матрицы}

Матрицы $A$ и $B$ называются \textbf{подобными}, если существует обратимая матрица $C$, такая что
\[
    B=C^{-1}AC.
\]
Подобные матрицы задают один и тот же оператор в разных базисах.

\section{Билет 26. Образ, ядро, ранг и дефект}

Пусть $L:V\to W$ --- линейный оператор.

\textbf{Образ оператора:}
\[
    \Image L=\{y\in W\mid \exists x\in V:\ y=L(x)\}.
\]

\textbf{Ядро оператора:}
\[
    \Ker L=\{x\in V\mid L(x)=0\}.
\]

$\Image L$ является подпространством в $W$, а $\Ker L$ --- подпространством в $V$.

\textbf{Ранг оператора:}
\[
    r(L)=\dim \Image L.
\]

\textbf{Дефект оператора:}
\[
    d(L)=\dim \Ker L.
\]

\subsection*{Теорема о ранге и дефекте}

Если $\dim V=n$, то
\[
    r(L)+d(L)=n.
\]

\subsection*{Идея доказательства}

В выбранном базисе условие $x\in\Ker L$ равносильно системе
\[
    AX=0,
\]
где $A$ --- матрица оператора. Пространство решений однородной системы имеет размерность $n-\rank A$. Так как $\rank A=r(L)$, получаем
\[
    d(L)=n-r(L),
\]
то есть
\[
    r(L)+d(L)=n.
\]

\section{Билет 27. Сумма и произведение операторов}

Пусть $\End_P V$ --- множество всех линейных операторов $V\to V$.

Для операторов задаются операции:
\[
    (L_1+L_2)(x)=L_1(x)+L_2(x),
\]
\[
    (\alpha L)(x)=\alpha L(x),
\]
\[
    (L_2L_1)(x)=L_2(L_1(x)).
\]

Относительно сложения и композиции $\End_P V$ является кольцом с единицей $I$.

Если $\dim V=n$ и выбран базис, каждому оператору $L$ соответствует его матрица $A_L$. Отображение
\[
    \varphi:\End_P V\to M_n(P),\qquad \varphi(L)=A_L
\]
является изоморфизмом колец:
\[
    A_{L_1+L_2}=A_{L_1}+A_{L_2},
\]
\[
    A_{\alpha L}=\alpha A_L,
\]
\[
    A_{L_2L_1}=A_{L_2}A_{L_1},
\]
\[
    A_I=E.
\]

\section{Билет 28. Алгебра над полем}

\textbf{Алгебра над полем $P$} --- это множество $A$, которое одновременно является кольцом и линейным пространством над $P$, причем умножение совместимо с умножением на скаляры:
\[
    \alpha(xy)=(\alpha x)y=x(\alpha y).
\]

Примеры:
\begin{itemize}
    \item $M_n(P)$;
    \item $\End_P V$;
    \item кольцо многочленов $P[x]$.
\end{itemize}

Если $\dim V=n$, то при фиксированном базисе
\[
    \End_P V\cong M_n(P)
\]
как алгебры над $P$.

Следствие:
\[
    \dim \End_P V=\dim M_n(P)=n^2.
\]

\section{Билет 29. Обратимые операторы и группа $GL$}

\textbf{Мультипликативная группа кольца} $R$ --- это множество всех обратимых элементов кольца относительно умножения. Обозначается $U(R)$.

Оператор $L\in\End_P V$ называется \textbf{обратимым}, если существует оператор $L^{-1}$, такой что
\[
    LL^{-1}=L^{-1}L=I.
\]

Множество всех обратимых операторов обозначается $\GL(V)$.

Множество всех обратимых матриц порядка $n$ над $P$ обозначается $\GL(n,P)$.

Так как оператор обратим тогда и только тогда, когда его матрица обратима, то
\[
    \GL(V)\cong \GL(n,P).
\]

\section{Билет 30. Теорема об обратимых операторах}

Пусть $V$ --- конечномерное пространство, $\dim V=n$, $L\in\End_P V$. Следующие условия эквивалентны:

\begin{enumerate}
    \item $L$ обратим;
    \item $L$ --- биекция;
    \item $\Image L=V$;
    \item $\Ker L=0$;
    \item для любого базиса $a_1,\ldots,a_n$ система $L(a_1),\ldots,L(a_n)$ тоже является базисом $V$;
    \item матрица $A_L$ обратима, то есть $\det A_L\ne0$.
\end{enumerate}

\subsection*{Короткая схема доказательства}

\begin{itemize}
    \item Обратимость равносильна биективности.
    \item Биективность дает $\Image L=V$ и $\Ker L=0$.
    \item Если $\Ker L=0$, то образы базисных векторов линейно независимы; так как их $n$, они образуют базис.
    \item Если образы базиса образуют базис, то каждый вектор имеет прообраз, значит оператор сюръективен, а в конечномерном случае обратим.
\end{itemize}

\section{Билет 31. Собственные векторы и собственные значения}

Пусть $L\in\End_P V$.

Ненулевой вектор $a$ называется \textbf{собственным вектором}, если существует $\lambda\in P$, такое что
\[
    L(a)=\lambda a.
\]

Число $\lambda$ называется \textbf{собственным значением}.

\textbf{Собственное подпространство}, соответствующее $\lambda$:
\[
    V_\lambda=\{a\in V\mid L(a)=\lambda a\}=\Ker(L-\lambda I).
\]

Нулевой вектор добавляют в это множество, но собственным вектором он не считается.

\subsection*{Примеры}

\begin{itemize}
    \item Для $L=\alpha I$ любой ненулевой вектор собственный, собственное значение равно $\alpha$.
    \item У проекции на ось $Ox$: векторы оси $Ox$ имеют $\lambda=1$, векторы оси $Oy$ имеют $\lambda=0$.
    \item Поворот плоскости на угол, не кратный $\pi$, не имеет вещественных собственных векторов.
\end{itemize}

\subsection*{Предложение}

Собственные векторы, отвечающие попарно различным собственным значениям, линейно независимы.

\subsection*{Идея доказательства}

Если
\[
    \alpha_1a_1+\cdots+\alpha_ka_k=0,
\]
применяем $L$, затем вычитаем первое равенство, умноженное на одно из собственных значений. Получаем линейную комбинацию меньшего числа собственных векторов. По индукции все коэффициенты равны нулю.

\section{Билет 32. Характеристический многочлен и спектр}

Пусть $A$ --- матрица оператора $L$ в некотором базисе.

\textbf{Характеристическая матрица:}
\[
    \lambda E-A.
\]

\textbf{Характеристический многочлен:}
\[
    f_L(\lambda)=\det(\lambda E-A).
\]

Корни характеристического многочлена называются \textbf{характеристическими корнями}.

\textbf{Спектр} $\sigma(L)$ --- множество характеристических корней с учетом кратностей.

\subsection*{Независимость от выбора базиса}

Если в другом базисе матрица оператора равна
\[
    B=C^{-1}AC,
\]
то
\[
    \det(\lambda E-B)
    =\det(C^{-1}(\lambda E-A)C)
    =\det(\lambda E-A).
\]
Значит, характеристический многочлен и спектр не зависят от выбора базиса.

\subsection*{Теорема}

Собственные значения оператора --- это в точности корни характеристического многочлена.

\subsection*{Короткое доказательство}

Если $L(a)=\lambda_0a$, $a\ne0$, то
\[
    (\lambda_0I-L)(a)=0.
\]
Значит, $\lambda_0I-L$ необратим, поэтому
\[
    \det(\lambda_0E-A)=0.
\]

Обратно, если $\det(\lambda_0E-A)=0$, то $\lambda_0I-L$ необратим, поэтому его ядро ненулевое. Значит, существует $a\ne0$, для которого
\[
    (\lambda_0I-L)(a)=0,
\]
то есть
\[
    L(a)=\lambda_0a.
\]

\section{Билет 33. Алгоритмы нахождения собственных значений и векторов}

\subsection*{Собственные значения}

\begin{enumerate}
    \item Выбрать базис и построить матрицу $A$.
    \item Составить характеристическую матрицу $\lambda E-A$.
    \item Вычислить характеристический многочлен
    \[
        f(\lambda)=\det(\lambda E-A).
    \]
    \item Найти корни $f(\lambda)$ в поле $P$.
\end{enumerate}

\subsection*{Собственные векторы для значения $\lambda_i$}

\begin{enumerate}
    \item Составить матрицу $\lambda_iE-A$.
    \item Решить однородную систему
    \[
        (\lambda_iE-A)X=0.
    \]
    \item Найти фундаментальную систему решений.
    \item Записать собственное подпространство $V_{\lambda_i}$ как линейную оболочку найденных базисных решений.
\end{enumerate}

\section{Билет 34. Диагонализируемые операторы}

Оператор $L$ называется \textbf{диагонализируемым}, если существует базис, в котором его матрица диагональна.

\subsection*{Критерий}

Оператор $L$ диагонализируем тогда и только тогда, когда в $V$ существует базис, состоящий из собственных векторов оператора $L$.

\subsection*{Доказательство}

Если матрица диагональна, то для базисных векторов $e_i$ имеем
\[
    L(e_i)=\lambda_i e_i.
\]
Значит, все $e_i$ --- собственные векторы.

Обратно, если есть базис из собственных векторов $e_i$, то
\[
    L(e_i)=\lambda_i e_i.
\]
Поэтому в этом базисе матрица оператора диагональна с числами $\lambda_i$ на диагонали.

\section{Билет 35. Критерий диагонализируемости через кратности}

Пусть характеристический многочлен оператора $L$ раскладывается на линейные множители:
\[
    f_L(\lambda)=\prod_i(\lambda-\lambda_i)^{k_i},
\]
где $\lambda_i$ --- различные собственные значения, а $k_i$ --- их алгебраические кратности.

Всегда верно:
\[
    \dim V_{\lambda_i}\le k_i.
\]

\subsection*{Критерий}

Оператор $L$ диагонализируем тогда и только тогда, когда для каждого собственного значения
\[
    \dim V_{\lambda_i}=k_i.
\]

Эквивалентно: сумма размерностей всех собственных подпространств равна $\dim V$.

\subsection*{Практический алгоритм}

Чтобы проверить диагонализируемость:
\begin{enumerate}
    \item найти характеристический многочлен;
    \item найти кратности его корней;
    \item для каждого корня найти $\dim\Ker(\lambda_iE-A)$;
    \item сравнить: если все геометрические кратности равны алгебраическим, оператор диагонализируем.
\end{enumerate}

\section{Билет 36. Инвариантные подпространства}

Подпространство $W\le V$ называется \textbf{инвариантным относительно $L$}, если
\[
    L(W)\le W,
\]
то есть для любого $x\in W$ имеем $L(x)\in W$.

\subsection*{Примеры}

\begin{itemize}
    \item $\Ker L$ инвариантно;
    \item $\Image L$ инвариантно для оператора $L:V\to V$;
    \item собственное подпространство $V_\lambda$ инвариантно;
    \item любое подпространство, содержащее $\Image L$, инвариантно.
\end{itemize}

\subsection*{Связь с собственными векторами}

Одномерное подпространство $\Lin{x}$ инвариантно относительно $L$ тогда и только тогда, когда $x$ --- собственный вектор.

\subsection*{Связь с блочно-треугольной матрицей}

Если $W$ --- инвариантное подпространство, $\dim W=k$, выбираем базис $e_1,\ldots,e_k$ в $W$ и дополняем до базиса всего $V$. Тогда матрица $L$ имеет вид
\[
    A=
    \begin{pmatrix}
        A_{11} & A_{12}\\
        0 & A_{22}
    \end{pmatrix}.
\]

Нулевая нижняя левая клетка появляется потому, что образы векторов из $W$ снова лежат в $W$.

\subsection*{Связь с блочно-диагональной матрицей}

Если
\[
    V=W\oplus U,
\]
где и $W$, и $U$ инвариантны относительно $L$, то в базисе, составленном из базисов $W$ и $U$, матрица имеет вид
\[
    A=
    \begin{pmatrix}
        A_{11} & 0\\
        0 & A_{22}
    \end{pmatrix}.
\]

То есть прямое разложение на инвариантные подпространства дает блочно-диагональную матрицу.

\section{Билет 37. Индуцированный оператор}

Пусть $W\le V$ --- инвариантное подпространство относительно $L$.

\textbf{Индуцированный оператор} $L_W:W\to W$ задается формулой
\[
    L_W(x)=L(x)
\]
для всех $x\in W$.

Он корректно определен, потому что $W$ инвариантно: если $x\in W$, то $L(x)\in W$.

\subsection*{Теорема}

Характеристический многочлен индуцированного оператора $L_W$ делит характеристический многочлен исходного оператора $L$.

\subsection*{Короткое доказательство}

Выбираем базис $e_1,\ldots,e_k$ в $W$ и дополняем его до базиса $V$. В этом базисе матрица $L$ имеет блочно-треугольный вид:
\[
    A=
    \begin{pmatrix}
        A_{11} & A_{12}\\
        0 & A_{22}
    \end{pmatrix},
\]
где $A_{11}$ --- матрица индуцированного оператора $L_W$.

Тогда
\[
    \det(\lambda E-A)
    =
    \det(\lambda E-A_{11})\det(\lambda E-A_{22}).
\]

Первый множитель $\det(\lambda E-A_{11})$ --- характеристический многочлен $L_W$, значит он делит характеристический многочлен $L$.

\subsection*{Следствие}

Если
\[
    V=W\oplus U,
\]
где оба подпространства инвариантны, то
\[
    f_L(\lambda)=f_{L_W}(\lambda)f_{L_U}(\lambda).
\]

\end{document}
